%
% 55-dimension.tex -- Homologie in höheren Dimensionen
%
% (c) 2026 Prof Dr Andreas Müller
%
\subsection{Homologie in anderen Dimensionen}
Die Homologiegruppe $H_1(U)$ von $U$ gibt Information über die 
Zusammenziehbarkeit von geschlossenen Kurven.
Damit lässt sich aber noch nicht entscheiden, ob $U$ aus verschiedenen,
nicht verbundenen Komponenten besteht.
Die Homologiegruppe $H_0$ löst dieses Problem.
Im dreidmensionalen Raum $U=\mathbb{R}^3\setminus\{0\}$ ohne den
Nullpunkt ist jede Kurve zusammeziehbar, die Homologiegruppe $H_1$
kann also das Fehlen des Nullpunktes nicht detektieren.
Dazu brauchte es die höhere Homologiegruppe $H_2(U)$.

%
% H_0: die Zusammenhangskomponenten von U
%
\subsubsection{$H_0(U)$: Die Zusammenhangskomponenten von $U$}
Eine zur Konstruktion der Homologiegruppe $H_1(U)$ analoge Konstruktion
kann auch eine Dimension tiefer durchgeführt werden.
In Definition~\ref{buch:integration:homologie:definition:randoperator}
wurden bereits die $0$-Ketten definiert und es wurde illustriert,
wie der Randoperator aus einer $1$-Kette eine $0$-Kette macht.
Die Menge
\[
B_0(U)
=
\{
\partial c
\mid
c\in C_1(U)
\}
\]
wird erzeugt von Paaren von Punkten, die durch eine Kurve 
verbunden sind.
Wir sagen wieder, zwei $0$-Ketten sind homolog, wenn sie sich
um den Rand einer $1$-Kette unterscheiden.
Damit lässt sich wieder eine Gruppe
\[
H_0(U)
=
C_0(U) / B_0(U)
=
c + B_0(U)
=
\{
c + \partial d
\mid
d\in C_1(U)
\}
\]
konstruieren.
Sie wird erzeugt von homologen Punkten.

\begin{satz}[Wegzusammenhang]
Wenn zu zwei beliebigen Punkten $P_1,P_2\in U$ die $0$-Kette
$P_2-P_1$ ein Rand einer $1$-Kette $c$ ist, dann gibt es einen Weg
zwischen $P_1$ und $P_2$.
Eine Teilmenge $U$ von $\mathbb{C}$ ist genau dann wegzusammenhängend,
\index{wegzusammenhangend@wegzusammenhängend}%
wenn es eine $1$-Kette $c \in C_1(U)$ gibt mit
$\partial c=P_2-P_1$.
\end{satz}

\begin{proof}
Wenn es einen Weg $\gamma$ zwischen $P_1$ und $P_2$ gibt, z.~B.~weil $U$
wegzusammenhängend ist, dann hat die
$1$-Kette $c=\gamma$ den verlangten Rand $\partial c = P_2-P_1$.

Wenn es eine $1$-Kette
\[
c
=
\sum_k
c_k\gamma_k
\]
mit $\partial c = P_2-P_1$ gibt, dann lassen sich die $\gamma_k$ zu einem
Weg zusammensetzen, der in $P_1$ beginnt und in $P_2$ endet.
Nicht unbeding alle Wege der Kette werden verwendet.
Einige Wege in der Kette bilden geschlossene Wege, die ignoriert
werden können.
Die Verkettung der Wege ist ein möglicherweise nicht differenzierbarer
Weg, somit liegen $P_1$ und $P_2$ in der gleichen Wegzusammenhangskomponente.
\end{proof}

Die Homologiegruppe $H_0(U)$ wird erzeugt von den
Zusammenhangskomponenten.
Natürlich enthält auch $H_1(U)$ Information über die
Zusammenhangskomponenten, da geschlossene Wege jeweils in einer
einzelnen Zusammenhangskomponente enthalten sind und daher
Wege in verschiedenen Zusammenhangskomponenten nicht homolog sein
können.

%
% Höhere Homologiegruppen
%
\subsubsection{Höhere Homologiegruppen}
Die Konstruktionen, die zur Homologiegruppe $H_1$ geführt haben,
lassen sich in jeder beliebigen Dimension durchführen.
Im folgenden sei daher $k>1$ eine natürliche Zahl.
Eine $k$-dimensionale Kette ist eine formale Linearkombination
von Abbildungen des $k$-dimensionalen Einheitswürfels
\[
\Gamma
\colon
I^k
\to
U
:
(t_1,\dots,t_k)
\mapsto
\Gamma(t_1,\dots,t_k).
\]
Die $(k-1)$-dimensionalen Seitenflächen des Würfels sind Abbildungen
der Form
\begin{align*}
\Gamma_i^0
&
\colon I^{k-1}\to U
:
(s_1,\dots,s_{k-1})
\mapsto
\Gamma(s_1,\dots,s_{i-1},0,s_{i},\dots,s_{k-1})
\intertext{und}
\Gamma_i^1
&
\colon I^{k-1}\to U
:
(s_1,\dots,s_{k-1})
\mapsto
\Gamma(s_1,\dots,s_{i-1},1,s_{i},\dots,s_{k-1})
\end{align*}
für $i=1,\dots,k-1$.
Der Randoperator ist wieder die lineare Abbildung, die auf $\Gamma$
als eine Linearkombination von $\Gamma^0_i$ und $\Gamma^1_i$ 
mit Koeffizienten $\pm1$.

Ein $k$-dimensionaler Zyklus ist eine $k$-dimensionale Kette $c$ mit
der Eigenschaft $\partial c=0$.
Die Menge der $k$-dimensionalen Zyklen wird mit
\[
Z_k
=
Z_k(U)
=
\{
c\in C_k(U)
\mid
\partial c=0
\}
\]
bezeichnet.

Der Randoperator auf $C_{k+1}$ liefert zu einem $(k+1)$-dimensionalen
Würfel die Oberfläche des Würfels, bestehend aus $k$-dimensionalen
Würfeln.
Die Menge der $k$-dimensionalen Ränder ist
\[
B_k
=
B_k(U)
=
\partial C_k(U).
\]
Zwei $k$-dimensionale Zyklen heissen homolog, wenn sie sich um
einen $k$-dimensionalen Rand unterscheiden, wir können daher auch
eine $k$-dimensionale Homologiegruppe
\[
H_k(U)
=
Z_k(U) / B_k(U)
=
\{
c + B_k(U)
\mid
c\in Z_k(U)
\}
\]
konstruieren.

%
% Homologiegruppen und Homotopie
%
\subsubsection{Homologiegruppen und Homotopie}
Man kann jeweils eines der Argument einer Abbildung $\Gamma$
eines $k$-dimensionalen Würfels als den Homotopieparameter
betrachten.
Die Abbildung wird dann als eine Homotopie zwischen zwei
$(k-1)$-dimensionalen Würfeln interpretiert.
Die Homologiegruppe betrachtet daher homotope $(k-1)$-dimensionale
Zyklen als homolog, wenn sie homotop sind.

Die Menge $U=\mathbb{R}^3$ ist zusammenziehbar, jeder Abbildung $\Gamma$
ist daher homotop zur konstanten Abbildung in den Nullpunkt.
Daher sind alle $k$-dimensionalen Zyklen auch Ränder und alle
Elemente in $Z_k(U)$ sind homolog.
Die Homologiegruppen sind daher $H_k(\mathbb{R})=\{0\}$.

%
% Homologiegruppen und Löcher
%
\subsubsection{Homologiegruppen und ``Löcher''}
Die Homologiegruppe $H_1(U)$ erkennt die Löcher in einem zweidimensionalen
Gebiet $U$.
Die Homologiegruppen $H_k$ mit $k>1$ erkennen höherdimensionale ``Löcher''.

Als Beispiel untersuchen wir die die Menge $U=\mathbb{R}\setminus \{0\}$.
Die Abbildung
\[
\Gamma
\colon
I^3
\to
\mathbb{R}^3
:
(t_1,t_2,t_3)
\mapsto
(
t_1-{\textstyle\frac12},
t_2-{\textstyle\frac12},
t_3-{\textstyle\frac12}
)
\]
ist keine Abbildung in $U$, da der Punkt $(\frac12,\frac12,\frac12)$ auf
den Nullpunkt abgebildet wird, der nicht in $U$ ist.
Der Randoperator liefert aber eine Linearkombination der Randquadrate,
die in $U$ enthalten sind.
Sie bilden einen Zyklus, da die Ränder der einzelnen Quadrate alle
in beiden Richtungen durchlaufen werden.
Dieser Zyklus ist also nicht in $B_2(U)$ enthalten.
Jede Abbildung von $I^3\to \mathbb{R}^3$, die den Nullpunkt trifft,
hat als Rand einen Zyklus in $U$, der nicht ein Rand sein kann.
Somit gibt es in der Homologiegruppe $H_2$ genau einen Zyklus, der
nicht homolog zur konstanten Abbildung ist.
Somit ist $H_2(U)=\mathbb{Z}$.
Die Homologiegruppe $H_2(U)$ hat also das Loch im Nullpunkt erkannt.
